\summarychapter{21}{A Song to Lady Poverty}{The Apostolic Lifestyle}

\begin{inbrief}
Against Innocent III's vision of Christ ruling the world through his vicar stood the older image of
the Savior who had nowhere to lay his head. In the twelfth and thirteenth centuries wandering
preachers of voluntary poverty---Arnold of Brescia, the Waldenses, the Cathari, Dominic and
Francis---challenged a wealthy church, some within its bounds and some outside them, and the church
answered
with preaching, crusade and Inquisition.
\question{Where is true Christianity found: in a sacramental institution or in a self-denying way of
life?}
\end{inbrief}

The troubadours of the twelfth century sang of courtly love. Francis of Assisi, ``Christ's very
courageous knight,'' sang of Lady Poverty. A story tells how he sought her in the mountains and found
her on ``the throne of her neediness.'' She had lived with Adam in Paradise, wandered homeless after
the fall, been chosen by Jesus and embraced by the monks, until Avarice made them rich and drove her
out. Francis took her as his bride. She had many other admirers. Growing numbers of laypeople read the
Bible in their own language and vowed to obey the word, ``Sell what you have, give to the poor, and
come, follow me.'' Some stayed orthodox and some became heretics, with at times only a knife edge
between the two.

\ssection{Poverty, Heresy, and Violence}
The poverty movement is a lasting reminder that political Christianity is only partial Christianity.
What good are canon law, crusades and disputations if ordinary people ask for bread and receive a
stone? The movement fed on resentment of a corrupt and neglectful priesthood and was often joined to
social restlessness, but its heart was spiritual hunger. Pastoral care was a lost art. Robert
Grosseteste, bishop of Lincoln (1235--1253), blamed the greed and immorality of the clergy on the
Roman court, which appointed ``not pastors but destroyers of men.'' Reformers since the tenth century
had called for a return to apostolic poverty, but such appeals had always been channeled into a new
monastic reform. Now some preachers would not stay within the church's bounds.

Moderns, who hold faith to be a private choice, find the medieval attitude to heresy hard to grasp.
In Christendom faith was the cement of society and its life-giving soul; to deny a doctrine was close
to treason, and heresy was no more tolerable than cancer in the body. Heresy meant a baptized
person's denial of any revealed truth, and since those truths included the church's unity and the pope's appointment by God, disobedience was itself heresy. The church aimed to convert the
heretic and to protect society. But how far may it go? Here it fell into its gravest inner conflict.
By creating the Inquisition, which not only executed heretics but tortured them, it chose a course
that cannot be reconciled with the kingdom it sought; in driving out one devil it let in seven. Few
saw the contradiction then. Almost everyone agreed that a pure church was God's will.

\ssection{The Poor in the Kingdom}
But where is the pure church: in the papal palace, or among the hungry who are fed? One of the first
to ask was Arnold, an abbot at Brescia in northern Italy, who preached that the vices of the clergy
came from the church's attempt to rule the world. It should hand its property and secular power to
the state and return to apostolic simplicity. Banished by Innocent II in 1139, he studied in Paris
under Abelard and was denounced by Bernard. After five years he appeared in Rome, where the people,
dreaming of the ancient republic, had seized power, and he became leader of the new secular
government, calling the College of Cardinals a den of thieves. After about ten years Hadrian IV put
Rome under interdict, Frederick Barbarossa captured Arnold, and in 1155 he was burned and his ashes
thrown into the Tiber.

Peter Waldo (c.~1140--1218), who had grown rich as a merchant in Lyons, was struck by a troubadour's ballad of Alexis,
who left his bride on their wedding night for the Holy Land and returned to die unrecognized as a
beggar in his parents' courtyard. A priest pointed Waldo to Jesus' word to the rich young ruler, the
text that had moved Anthony. He provided for his wife, placed his daughters in a cloister and gave
the rest to the poor. He had two priests translate parts of the Bible into French, memorized long
passages and taught common people to imitate Christ. His innovation was to apply poverty,
discipleship and preaching to all true Christians, not only monks. He sent followers out two by two;
they called themselves the ``Poor in Spirit,'' and we call them Waldenses. Forbidden by the archbishop
of Lyons to preach, Waldo answered, ``We must obey God rather than men,'' and was excommunicated. At
the Third Lateran Council (1179) the Waldenses were mocked by the Englishman Walter Map. Pope
Alexander III found no heresy in them but allowed them to preach only at the invitation of bishops.
They preached anyway, spread into southern France and Italy, and in 1184 Lucius III excommunicated
them.

Their aim was a purer church, recovered by living as simply as the apostles had, which meant giving up
worldly power; Rome could not renounce its sacraments, its priesthood or its authority. The
Waldenses increasingly held that only Christ's teaching binds and Scripture must rule. To survive
they accepted two levels of commitment, the vowed ``Poor'' and a circle of ``friends'' who remained
Catholics and supplied recruits and support. Evangelicals have often called them reformers before the
Reformation. Their appeal to the Bible does sound like Luther or Calvin, but they pictured salvation as a penitent and poor life, without the Reformation's clear note of grace.

\ssection{The Peril of New Heresies}
Most troublesome were the Cathari (``pure ones''), also called Albigenses after the town of Albi in
southern France. What is known of them comes mostly from their enemies. They probably entered Europe
from Bulgaria, where a branch was called Bogomiles. Like the Gnostics they taught an eternal conflict
of a good and an evil power. Matter and the body are the work of the evil power, the god of the Old
Testament, who imprisoned the soul. The true Cathar avoided marriage, sexual intercourse, meat and
possessions: a radical poverty, but rooted in the nature of the universe rather than the example of
Jesus. Christ was not human but a life-giving spirit who taught the way, like a Buddha; salvation by
his death on the cross was unthinkable. Arnold and Waldo had refused obedience; the Cathari rejected
basic Christianity, seeking escape from evil by liberating the pure self rather than by repentance
and faith. By 1200 they enjoyed the protection of the princes of Toulouse and were spreading fast.
The church had three weapons: preaching, crusade and Inquisition.

\ssection{Rise of the Dominicans}
The pope's preachers failed until the Spaniard Dominic Guzman (1170--1221) saw why in 1206. They
relied on pomp, which the Albigenses took as a mark of false religion. Dominic went among them poor,
barefoot and begging. His peaceful mission lasted only two years before Innocent III turned to
force, but he gathered like-minded men and went on elsewhere, and in 1220 the order was approved. The
Dominicans were mendicants (beggars) and friars (brothers), who, unlike monks, lived among the people
to preach and teach; as monasteries had served the countryside, the friars served the towns. Innocent
meanwhile called a crusade against Christian heretics. The northern French, eager to seize the
separate country of the south, came pillaging and murdering with a brutality that shocked even the
pope. By 1215 Toulouse had been cleared of Albigenses and northern lords held the ruined south.

\ssection{Enter the Inquisitor}
The scattered heretics had to be hunted out, and that was the work of the Inquisition, a name that
still stands for the miscarriage of justice. It began in 1184, when Lucius III required bishops to
``inquire'' into their people's beliefs. In 1215 the Fourth Lateran Council under Innocent III
provided for punishment of heretics by the state, confiscation of their property, excommunication of
those who would not act and full forgiveness for those who cooperated. In 1220 the pope transferred
the Inquisition from the bishops to the Dominicans, and in 1229 the Synod of Toulouse fixed a
procedure that left the accused almost no rights. The inquisitor answered only to the pope and was
prosecutor and judge; the trial was secret; the accused had to prove his innocence without counsel
and without knowing his accusers. In 1252 Innocent IV authorized torture, which earlier popes and
theologians had rejected with horror. Because canon law forbade a cleric to shed blood, the condemned
were handed to the civil power, usually to be burned. Almost everyone since Augustine agreed that a
rotten limb must be cut off to save the body.

Crusade and Inquisition ended Catharism before the century was out. The Waldenses, whom the Inquisition also hunted, held out in the Italian Alps until the Reformation, which they greeted as an ally. The
Inquisition's weakness was that it could amputate but not heal. Healing came from Assisi, north of
Rome. Francis (1182--1226), born Giovanni Bernardone, son of a rich cloth merchant, sought military
glory until capture, illness and visions turned him around. He gave up wealth and honor, embraced
the lepers who had once repelled him, and in a ragged cloak with a rope belt wandered with a few
companions, asking alms of the wealthy, passing them on to the needy and praising the happiness of a life as poor as the apostles'.

\ssection{Founding the Franciscans}
In 1209 Francis drew up a simple Rule made up mostly of sayings of Christ (the call to carry one's cross, the counsel given to the rich young ruler, the orders for the apostles' mission), and took it to Innocent III. The
scene repeated Waldo's in 1179, but Innocent saw his predecessors' mistake and approved the little
band. Francis called them Friars Minor, ``lesser brothers''; we call them Franciscans. His vision was
for the world. After failed attempts to reach Syria and Morocco he went with a crusading army to
Egypt in 1219, tried to convert the sultan and visited the holy places. In his absence the brothers
divided between those who wanted more organization and those who clung to the original poverty.
Francis was a model, not a manager. He asked for Cardinal Ugolino as adviser and handed the
administration to Peter de Cataneo. Ugolino, later Pope Gregory IX, saw the movement as an instrument
for restoring Roman authority where Arnold, Waldo and the Cathari had undermined it; Francis wanted
to reform the world by Christlike humility. In 1223 Honorius III confirmed a new Rule that allowed
an elaborate organization and made begging fundamental. Franciscans thus joined Dominicans in a great
mendicant counteroffensive. Francis died at Assisi on October 3, 1226, singing.

Herbert Workman wrote that the Sermon on the Mount was, briefly, actually lived; then the
dream faded for all but a few. For another century the order was split between the Conventuals, who
let the church hold property for their use, and the Spirituals, who held complete poverty to be a
mark of the true church. Lady Poverty went on rebuking a rich and powerful church, but like Innocent's universal dominion she proved impossible for the whole church for all time.

\begin{keyfacts}
\keypeople{Francis of Assisi; Robert Grosseteste; Arnold of Brescia; Peter Waldo; Alexander III;
Lucius III; Dominic Guzman; Innocent III; Innocent IV; Cardinal Ugolino (Gregory IX); Honorius III.}
\keydates{%
  \kd{1155}{Arnold of Brescia burned at Rome}
  \kd{1179}{Waldenses appeal to the Third Lateran Council}
  \kd{1184}{Lucius III excommunicates the Waldenses and orders bishops to ``inquire'' into heresy}
  \kd{1209}{Francis of Assisi wins approval for his first Rule}
  \kd{1215}{Fourth Lateran Council provides for the punishment of heretics; Albigenses crushed in Toulouse}
  \kd{1220}{Dominican order approved; Inquisition entrusted to the Dominicans}
  \kd{1223}{Honorius III confirms the new Franciscan Rule}
  \kd{1226}{Death of Francis of Assisi (October 3)}
  \kd{1252}{Innocent IV authorizes torture in the Inquisition}}
\keyterms{Apostolic poverty; heresy; Waldenses (``Poor in Spirit''); Cathari or Albigenses;
Bogomiles; dualism; mendicant; friar; Dominicans; Franciscans (Friars Minor); Inquisition;
Conventuals and Spirituals.}
\end{keyfacts}
